This section complements the model section of the main text and the ODD description above with the details that
the main text leaves out: the refuge capacity and the update scheme, the dimensionless groups and time scales, and
the candidates and rules that are not used in the main text.

\subsection{Hard capacity of the refuges and update scheme}
\label{sm:model-capacity}

The admission rule of the refuges, iterated to a fixed point, is given in the ODD description above. A single-pass
admission that does not count returning animals and newborns, so that a refuge can hold up to 11 individuals, is a
robustness variant; it leaves the primary outcome unchanged (0.995 with either rule).

Both the preference and the hard capacity are necessary with synchronous moves. Without a preference ($\pi=0$) the
refuges stay empty, because a deteriorated individual surrounded by acquaintances prefers an occupied neighboring
cell; with a soft capacity, simultaneous moves produce stampedes of deteriorated animals into the same refuge. With a
random sequential update a soft capacity also produces a persistent isolated class (O17p 1.00 and 0.995 with hard and
soft capacity), with a higher peak density ($\rho_{\rm pk}=0.43$ against 0.36). The hard capacity is therefore a
remedy for the collisions of the synchronous update, not a mechanism of its own. The total refuge capacity,
$f_Rc_R|G|=540$ places, is about 20\% of the adults at the peak, $\phi_R=f_Rc_R|G|/N_{\rm ad}(t_{\rm pk})\approx0.2$,
and bounds the persistently isolated fraction, $\Phi_{\rm BO}^{\rm pers}\le\phi_R$. The bound holds at all 50 points
of the refuge plane, so a class of $\ge10\%$ needs $f_Rc_R\gtrsim0.3$; in that plane a persistent class appears only
for $f_Rc_R\ge0.4$ and a strong preference.

\subsection{Dimensionless groups and time scales}
\label{sm:model-scales}

The values are for \textsf{C1}; Table~\ref{tab:sm-model-scales} compares the scales with Universe~25.
\begin{itemize}
  \item $\Pi_\gamma=\gamma/r=1$: crowding damage per excess neighbor relative to the weekly learning rate.
  \item $\Lambda=r\,a_w=1.2$: the largest competence a juvenile can learn inside the window. A socialized lineage
    sustains itself only if $w+\Lambda\langle\bar s_V\rangle\gtrsim1$; fast growth lowers $\langle\bar s_V\rangle$,
    because novices outnumber teachers (a \emph{learning dilution}). When $a_w$ is varied as well, $r$ and $a_w$ do not
    act through their product alone: the local weight of $a_w$ relative to $r$ is about one half between $a_w=6$ and
    12 weeks and about one between 12 and 24.
  \item $n^*/m_c\approx0.8$, with $n^*=1/\kappa-1/\alpha$ the preferred group size of a deteriorated individual.
    Without refuges it decides between pools, isolation, or both; with refuges it only marks the onset of
    homogenization, near $n^*\approx3$--4.
  \item $\phi_R\approx0.2$, the refuge capacity per adult at the peak (Sec.~\ref{sm:model-capacity}).
  \item $n_0=N_0/|G|=0.44$ and the Damk\"ohler number $\mathrm{Da}=t_{\rm mix}/t_{\rm det}$. Here
    $t_{\rm mix}=L^2/(4D_{\rm eff})$, with $D_{\rm eff}\approx0.13$ cells$^2$ per week the mean squared displacement
    over 16 weeks divided by $4\times16$ (8 exploratory runs), so $t_{\rm mix}\approx1.7\times10^3$ weeks; and
    $t_{\rm det}=t_{1/2}-t_{\rm on}\approx10$ weeks is the time the animals past the window (age $\ge a_w$) need to
    lose half of their competence ($t_{1/2}=16$, the first week with their mean $s<0.5$) after crowding sets in
    ($t_{\rm on}=6$, the first week with 5\% of the individuals in cells with $m>m_c$). These animals no longer learn,
    and newborns do not enter their mean until they leave the window, so $t_{\rm det}$ measures deterioration and not
    the dilution by pups born with $s\le0.2$. Thus $\mathrm{Da}\approx1.7\times10^2$, and deterioration is correlated
    over $\xi=(4D_{\rm eff}t_{\rm det})^{1/2}\approx2$ cells, the size of the contact neighborhood. In founder colonies
    the relevant transport is the colonization front, and $t_{\rm fill}\approx t_{\rm pk}\approx170$--215 weeks against
    $t_{\rm det}\approx20$--30 gives $\mathrm{Da}_{\rm fill}\approx10$. In Universe~25 mice crossed the enclosure
    daily and the social organization took a whole phase~C, about 35 weeks, to die, so
    $\mathrm{Da}\lesssim4\times10^{-3}$: the spatial structure of the model emerges among animals that are almost
    sessile on the time scale of deterioration, which is not Calhoun's regime. Colonies grow from local nuclei, so the
    peak density falls with $n_0$, from 0.40 at $N_0=400$ to 0.03 for eight founders.
  \item $\Gamma=t_{\rm det}/T_{\rm gen}\approx1.1$, with $T_{\rm gen}\approx9$ weeks the interval between the median
    birth weeks of the first and second generations; the interquartile range over runs is 1.0--1.29, and the
    founders, a closed cohort, give 1.2. The mean over all adults (age $\ge a_{\min}$) would give 0.8
    because it is diluted by the first generation, which enters it with $s\le0.2$ while still inside the window. With
    overlapping cohorts the interval between median birth weeks is not a demographic generation time; in the null
    models it drops to 5 weeks, below the age of first reproduction, which inflates their $\Gamma$ (1.4--1.6). With the
    mean age of the mothers at birth ($T_{\rm gen}=14.3$ weeks in \textsf{C1}) $\Gamma\approx0.7$, and 0.60--0.85 in
    the other arms where it was computed: $\Gamma$ is of order one under either definition. It does not set the depth
    of the boom: with five times less crowding damage ($\gamma=0.02$) $\Gamma=1.6$ and the boom is still one to two
    generations deep (mean generation of the births 1.42 against 1.35, net reproductive number of the first generation
    0.52, O17p in 17 of 20 runs), whereas pups born at $s_{\rm nb}\approx0.4$ ($s_0=0.25$) deepen it more at about the
    same $\Gamma$ (1.3; mean generation 1.51, net reproductive number 0.62). In Universe~25 the deterioration took the
    whole phase~C, about 35 weeks, with $T_{\rm gen}\approx10$ weeks, so $\Gamma\approx3.5$ (our estimate, with a
    different operational definition): an order-of-magnitude comparison, not a test.
  \item \emph{Extinction time.} Without mortality before senescence $T_{\rm ext}=\max_i(b_i+\ell_i)\simeq
    t_{\rm last}+\ell_{\max}$, the last birth plus the longest residual lifetime. In \textsf{C1} $\ell_{\max}\approx110$
    weeks, with an interquartile range of a few weeks, of the order of $1/k\approx7$ weeks, so the spread of
    $T_{\rm ext}$ is that of $t_{\rm last}$, set by how abruptly natality stops,
    $\tau_B=(t_{B99}-t_{B90})/\ln10\approx7$ weeks. At fixed density the median of $T_{\rm ext}$ grows with the
    logarithm of the system size, from 162 weeks at $L=15$ to 184 at $L=60$, and its interquartile range over the
    median is 0.15 for $L\le20$ and about 0.08 for $L\ge30$. Ratios normalized by the peak time depend on where it falls
    on the plateau: $(T_{\rm ext}-t_{\rm pk})/t_{\rm pk}\approx2.1$ (Universe~25: $\ge1.84$) becomes 4.0 with the start
    of the plateau as peak time, and $t_{\rm pk}/a_{\max}\approx0.45$ (0.28), against 0.7 in Universe~25.
\end{itemize}

\begin{table}[t]
  \caption{Scales of \textsf{C1} and of Universe~25~\cite{Calhoun1973}. $^\dagger$Our estimate. The model's
  $t_{\rm det}$ is measured on the animals past the window; for Universe~25 it is the duration of phase~C, so the two
  values of $\Gamma$ compare only in order of magnitude. Neither $|G|$ nor $m_c$ was calibrated to Universe~25, so
  densities are compared only in sign.}
  \label{tab:sm-model-scales}
  \begin{ruledtabular}
  \begin{tabular}{lll}
    Quantity & Model (\textsf{C1}) & Universe~25 \\
    \hline
    Resting sites & 900 cells & 256 nest boxes \\
    Comfortable occupancy & $m_c=8$ & 15 per box \\
    $K_\gamma$ (sites $\times$ occupancy) & 7200 & 3840 \\
    Isolation sites & 270 refuges, 540 places & high boxes, 25--50\%$^\dagger$ \\
    Peak; $N_{\rm pk}/K_\gamma$ & 2600; 0.36 & 2200; 0.57 \\
    Empty sites at the peak & 16\% & 20\% \\
    Movement & $D_{\rm eff}\approx0.13$ cells$^2$/week & enclosure crossed daily \\
    Maturity / menopause / life & 6 / 80 / 120 weeks & 6--7 / 80 / 114 weeks \\
    Generation time $T_{\rm gen}$ & $\approx9$ weeks & $\approx10$ weeks$^\dagger$ \\
    Deterioration time $t_{\rm det}$ & 10 weeks & $\approx35$ weeks (phase C)$^\dagger$ \\
    $\mathrm{Da}=t_{\rm mix}/t_{\rm det}$ & $\approx1.7\times10^2$ & $\lesssim4\times10^{-3}$ \\
    $\Gamma=t_{\rm det}/T_{\rm gen}$ & 1.1 & $\approx3.5^\dagger$ \\
  \end{tabular}
  \end{ruledtabular}
\end{table}

\subsection{Other candidates and rejected rules}
\label{sm:model-other}

\emph{Mortality.} The full form is
$p_{\rm death}(a,s)=\bigl[1+e^{-k(a-a_{\max}[1-\mu_s(1-s)])}\bigr]^{-1}$, where $\mu_s\ge0$ would advance the
senescence of deteriorated individuals; every model, candidates and controls, has $\mu_s=0$, so adult death is purely
senescent. There is no gestation
delay, and a female may give birth every week.

\textsf{C0} is the core model with the long-lived demography and R1, R2 and AV, without refuges
(\texttt{calhoun\_C0}). It is the minimal model that satisfies the essential patterns evaluable with $N_0=400$ when
O13 is read instantaneously (O17 in 0.92 of the runs), but its isolated class is transient (O17p in 0). In
\textsf{C0} each of R1, R2 and AV is necessary: without R2, O11 and O13 are never reproduced; without AV the
deteriorated animals collapse into a few pools, 76\% of the lattice is empty and the isolated class disappears; without
R1 adults recover when the colony thins and O17 falls to 0.40. \textsf{C2} (\texttt{calhoun\_C2}) is \textsf{C0} plus
R7 with the retreat driven by the competence of the mother at birth, $b_i=h_i$, with $q=1$ and $\pi=1000$: pups of
deteriorated mothers retreat, and socialized adults that later deteriorate do not. It is motivated by the beautiful
ones belonging to ``the last half of the population born''~\cite{Calhoun1973}; it makes the isolated class younger
than the crowded one, but the late-cohort pattern O13b only reaches MARGINAL, and O17p falls to 0.705 [0.64, 0.76].
Basing the retreat on juvenile competence (the value of $s$ at age $a_w$) fails, because under R2 almost every cohort
born after the first ends the window with $s<0.1$.

Several rules were unnecessary or destroyed some pattern, and are not part of any candidate: long-range jumps, an
attraction that decreases with $s$ (it homogenizes the population and raises the peak), territoriality and defeat of
males, social mortality, bonds restricted to competent pairs, and movement sub-steps (with more sub-steps the
coexistence of classes is lost).
